% Three-panel \method overview, deployment-time action improvement, and task.
\ifdefined\PaperPaletteLoaded\else% Central colour palette shared by the paper and all active figures.
% Method colours identify methods, not dataset conditions or phases.
\def\PaperPaletteLoaded{1}

% Paper-wide method and neutral colours.
\definecolor{paperCorePivot}{HTML}{008577}
\definecolor{paperCorePivotLight}{HTML}{E5F3EF}
\definecolor{paperIQL}{HTML}{3F6C8E}
\definecolor{paperTD3BC}{HTML}{8A681F}
\definecolor{paperCalQL}{HTML}{397C70}
\definecolor{paperBC}{HTML}{7654A3}
\definecolor{paperBCLight}{HTML}{F0EBF7}
\definecolor{paperPD}{HTML}{303030}
\definecolor{paperNeutralLight}{HTML}{F7F8F9}
\definecolor{paperReference}{HTML}{183B65}
\definecolor{paperInk}{HTML}{202B35}
\definecolor{paperGray}{HTML}{7A8286}
\definecolor{paperBackground}{HTML}{F9FAFB}
\definecolor{paperGrid}{HTML}{BFC3C6}
\definecolor{paperSummary}{HTML}{565656}
\definecolor{paperBand}{HTML}{777777}
\definecolor{paperSilver}{HTML}{AEB4BC}
\definecolor{paperSilverLight}{HTML}{EDF0F3}
\definecolor{paperSilverEdge}{HTML}{626A73}
\definecolor{paperPropellerLight}{HTML}{E4EEF7}
\definecolor{paperPropellerEdge}{HTML}{3974A1}

% Semantic aliases used by the framework overview.
\colorlet{hmBehavior}{paperBC}%
\colorlet{hmInk}{paperInk}%
\colorlet{hmGray}{paperGray}%
\colorlet{hmPaper}{paperBackground}%
\colorlet{overviewPolicyEdge}{paperCorePivot}%
\colorlet{overviewPolicyFill}{paperCorePivot!25}%
\colorlet{overviewActionEdge}{paperInk!70}%
\colorlet{overviewActionOuterFill}{paperGray!20}%
\colorlet{overviewDPEdge}{paperInk}%
\colorlet{overviewDPFill}{paperPropellerLight}%
\colorlet{droneAccent}{paperSilver}%
\colorlet{droneLight}{paperSilverLight}%

% Dataset-construction figure colours.
\definecolor{hdBlue}{RGB}{0,114,178}
\definecolor{hdPurple}{RGB}{145,94,155}
\definecolor{hdGood}{RGB}{0,130,115}
\definecolor{hdMedium}{RGB}{216,155,38}
\definecolor{hdBad}{RGB}{185,80,77}
\definecolor{hdInk}{RGB}{45,52,64}
\definecolor{hdPaper}{RGB}{248,250,252}

% Legacy illustration names retained for compatibility with optional figures.
\definecolor{policyedge}{RGB}{203,55,111}
\definecolor{policycoreedge}{RGB}{215,50,32}
\definecolor{policyfill}{RGB}{233,198,213}
\definecolor{policycorefill}{RGB}{236,196,199}
\definecolor{actionedge}{RGB}{61,77,150}
\definecolor{actionfill}{RGB}{159,154,199}
\definecolor{actionouterfill}{RGB}{159,208,240}
\definecolor{actoredge}{RGB}{155,77,152}
\definecolor{dpedge}{RGB}{7,10,28}
\definecolor{dpfill}{RGB}{227,228,233}
\definecolor{hmBlue}{RGB}{0,114,178}%

% Shared critic-landscape colour map.
\pgfplotsset{
  colormap={referencecriticmap}{
    color(0cm)=(paperCorePivot!5)
    color(1cm)=(paperCorePivot!22)
    color(2cm)=(paperCorePivot!65)
  }
}
\fi%
\begingroup
%
\begin{tikzpicture}[
  x=1cm,
  y=0.90cm,
  >=Stealth,
  font=\sffamily\scriptsize,
  panel/.style={rounded corners=4pt,draw=none,fill=none,
    line width=1.15pt},
  criticdetail/.style={panel,draw=none,fill=none},
  taskdetail/.style={panel,draw=none,fill=none},
  box/.style={rounded corners=2.5pt,draw=black,fill=white,
    line width=0.70pt,align=center,inner sep=4pt},
  databox/.style={box,draw=black,fill=paperNeutralLight},
  corebox/.style={box,draw=black,line width=0.82pt,fill=paperCorePivotLight},
  behaviorbox/.style={box,draw=black,densely dashed,fill=paperBCLight},
  plantbox/.style={box,draw=black,fill=overviewDPFill},
  paneltitle/.style={font=\sffamily\bfseries\small,text=hmInk,
    anchor=west},
  stage/.style={font=\sffamily\bfseries\scriptsize,text=hmInk!84,
    align=center},
  annotation/.style={draw=none,fill=none,inner sep=1pt},
  note/.style={annotation,font=\sffamily\fontsize{6.0}{6.7}\selectfont,
    text=hmInk!72,align=center},
  flow/.style={-{Stealth[length=5pt,width=5.4pt]},draw=hmInk,
    line width=0.78pt,rounded corners=2pt},
  baseline/.style={-{Stealth[length=4.2pt,width=4.6pt]},draw=hmInk!45,
    line width=0.50pt,densely dashed,rounded corners=2pt},
  criticflow/.style={-{Stealth[length=5pt,width=5.4pt]},draw=paperCorePivot,
    line width=0.95pt,rounded corners=2pt},
  supportflow/.style={-{Stealth[length=5pt,width=5.4pt]},draw=hmBehavior,
    densely dashed,line width=0.90pt,rounded corners=2pt},
  mainflow/.style={-{Stealth[length=5.3pt,width=5.6pt]},draw=paperCorePivot,
    line width=1.20pt,rounded corners=2pt},
  feedback/.style={-{Stealth[length=4.5pt,width=4.9pt]},draw=hmInk,
    densely dotted,line width=0.90pt,rounded corners=2pt,
    preaction={draw=white,line width=1.8pt}},
  axis/.style={-{Stealth[length=4.2pt,width=4.7pt]},draw=black!65,
    line width=0.68pt},
  actorpath/.style={-{Stealth[length=4.8pt,width=5.2pt]},draw=paperCorePivot,
    line width=1.00pt,rounded corners=2pt},
  behaviorpath/.style={-{Stealth[length=4.8pt,width=5.2pt]},draw=hmBehavior,
    densely dashed,
    line width=1.00pt,rounded corners=2pt}
]

% ===========================================================================
% (a) Explicit offline and online phases.
% ===========================================================================
% Reflowed layout: explicit positions and content-sized blocks, no stretching.
% Lift the lower model row; preserve node sizes and every arrow connection.
\ifdefined\CompactMethodPanels\else
\node[panel,
  minimum width=16.00cm,minimum height=5.895cm]
  (overviewPanel) at (8.00,8.225) {};

% Subtle phase regions, separated by a deployment boundary.
\path[rounded corners=3pt,fill=white,draw=none]
  (0.18,5.13) rectangle (8.10,10.28);
\path[rounded corners=3pt,fill=white,draw=none]
  (8.40,5.13) rectangle (15.82,10.28);
\draw[black!45,densely dashed,line width=0.72pt]
  (8.25,5.23) -- (8.25,10.30);

\ifdefined\HideInsetTitles\else
\node[paneltitle,text=black]
  at (0.22,11.23) {(a) \method{} framework overview};
\fi
\node[font=\sffamily\bfseries\small,text=paperCorePivot!80!black]
  at (4.10,10.78) {OFFLINE PHASE};
\node[note,font=\sffamily\fontsize{5.8}{6.3}\selectfont]
  at (4.10,10.43) {fit from recorded behaviour; regularise conservatively};
\node[font=\sffamily\bfseries\small,text=paperCorePivot!80!black]
  at (12.10,10.78) {DEPLOYMENT PHASE};
\node[note,font=\sffamily\fontsize{5.8}{6.3}\selectfont]
  at (12.10,10.43) {frozen models; repeat at each control step};

% ---------------------------------------------------------------------------
% Offline phase: fixed data, preprocessing, and model learning.
% ---------------------------------------------------------------------------
\node[databox,draw=black,line width=0.82pt,
  minimum width=2.10cm,minimum height=1.45cm,
  text width=1.84cm,font=\sffamily\fontsize{7.0}{8.2}\selectfont]
  (dataset) at (1.25,8.30)
  {\textbf{Offline dataset $\mathcal D$}\\[2pt]
   {\fontsize{6.2}{7.2}\selectfont $(o,a,r,o',a'_{\mathcal D},d)$}};
\node[note,font=\sffamily\fontsize{6.2}{7.2}\selectfont,
  text width=2.05cm] at (1.25,6.85)
  {fixed data ---\\no online exploration};
\node[annotation,font=\sffamily\fontsize{6.2}{7.2}\selectfont,
  text=hmInk!76]
  at (1.25,7.30) {standardise features};

% Extend the dataset hand-off and use its endpoint as the common split point.
\coordinate (offlineBranch) at (2.75,8.30);
\draw[flow,draw=black,line width=0.84pt]
  (dataset.east) -- (offlineBranch);

\node[corebox,minimum width=2.10cm,minimum height=1.25cm,
  text width=1.89cm,inner sep=3pt,
  font=\sffamily\fontsize{6.7}{8.0}\selectfont]
  (calqlTrain) at (4.05,9.20)
  {\textbf{Conservative\\behaviour-return fitting}\\[3pt]SARSA + $\mathcal R_{\mathrm{CQL}}$};

\node[corebox,minimum width=2.60cm,minimum height=1.25cm,
  text width=2.37cm,inner sep=3pt,
  font=\sffamily\fontsize{6.8}{8.0}\selectfont]
  (criticModel) at (6.85,9.20)
  {\textbf{\mbox{Conservative critic}}\\[2pt]
   Learned value landscape\\[2pt]frozen after fitting};


\node[behaviorbox,minimum width=3.00cm,minimum height=1.20cm,
  text width=2.72cm,inner sep=4pt,
  font=\sffamily\fontsize{7.0}{8.2}\selectfont]
  (behavior) at (4.65,6.45)
  {\textbf{Conditional behavior\\model}\\[3pt]
   fit once; freeze};

\draw[flow] (offlineBranch) |- (calqlTrain.west);
\draw[supportflow] (offlineBranch) |- (behavior.west);

% BC proposals regularise the recorded-behaviour return fit.
\draw[supportflow] (behavior.north)--(4.65,7.50)-|(calqlTrain.south);
\node[note,text width=2.1cm] at (5.90,7.55) {BC + uniform\\conservative samples};
\draw[criticflow] (calqlTrain.east) -- (criticModel.west);

% Explicit offline-to-online deployment interface.
\node[box,draw=black,fill=white,line width=0.72pt,
  rounded corners=2pt,inner sep=4pt,
  font=\sffamily\fontsize{6.7}{8.0}\selectfont,
  text=hmInk!82,align=center,minimum width=1.85cm,
  minimum height=1.95cm,text width=1.59cm]
  (frozenModels) at (8.10,6.50)
  {\textbf{Frozen models}\\[1pt]
   at deployment\\[4pt]
   \textcolor{paperCorePivot!80!black}{Critics}\\[3pt]
   \textcolor{hmBehavior}{BC model}};

% Drop below the critic, turn right, then turn down into the frozen models.
\draw[criticflow,line width=0.74pt]
  (criticModel.south) -- ++(0,-0.55) -| (frozenModels.north);
\draw[supportflow,line width=0.72pt]
  (behavior.east) -- (frozenModels.west |- behavior.center);

% ---------------------------------------------------------------------------
% Online phase: observe, optimize, act, and repeat in closed loop.
% ---------------------------------------------------------------------------
\node[databox,minimum width=0.74cm,minimum height=0.52cm]
  (observation) at (8.70,8.15) {$o_t$};

\node[corebox,draw=black,line width=1.08pt,
  minimum width=2.40cm,minimum height=1.85cm,text width=2.14cm,
  inner sep=4pt,font=\sffamily\fontsize{7.0}{8.3}\selectfont]
  (\method) at (10.85,8.15)
  {\textbf{\method}\\[1pt]
   Sample candidates\\[3pt]
   Rank by critic\\[3pt]
   Refine the best\\[3pt]
   Apply $u_t=u_{\max}a^\star$};

\node[circle,draw=black,fill=paperCorePivot!15,line width=1.18pt,
  minimum size=0.76cm,inner sep=1pt,font=\sffamily\bfseries\small]
  (optimalAction) at (12.95,8.15) {$u_t$};
\node[plantbox,minimum width=1.65cm,minimum height=1.65cm]
  (overviewPlant) at (14.65,8.15) {};
\node[annotation,font=\sffamily\bfseries\fontsize{6.7}{7.8}\selectfont,
  text=black,align=center]
  at ([yshift=-0.25cm]overviewPlant.north) {Drone task};

% Compact reference path and quadrotor glyph inside the plant block.
\draw[paperReference,densely dashed,line width=0.46pt]
  (13.85,7.82) .. controls (14.12,8.54) and (14.80,7.66) ..
  (15.40,8.45);
\begin{scope}[shift={(14.65,8.12)},rotate=16,scale=1.30]
  \draw[paperSilverEdge,line width=0.82pt] (-0.30,-0.19) -- (0.30,0.19);
  \draw[paperSilverEdge,line width=0.82pt] (-0.30,0.19) -- (0.30,-0.19);
  \shade[inner color=paperSilverLight,outer color=paperSilver,
    draw=paperSilverEdge] (0,0) circle[radius=0.088];
  \foreach \x/\y in {-0.30/-0.19,0.30/0.19,-0.30/0.19,0.30/-0.19} {
    \draw[paperPropellerEdge,fill=paperPropellerLight,line width=0.55pt]
      (\x,\y) ellipse[x radius=0.135,y radius=0.062];
  }
\end{scope}
\node[note,text=hmInk!78,font=\sffamily\fontsize{5.3}{6.0}\selectfont]
  at (14.65,7.47) {disturbed tracking};

% One straight forward signal path: observation, controller, action, and plant.
\draw[flow] (observation.east) -- (\method.west);
\draw[mainflow] (\method.east) -- (optimalAction.west);
\draw[mainflow] (optimalAction.east) -- (overviewPlant.west);

% The boundary bundle collects the learned outputs without long cross-panel
% routes. Three short local connections expose their online roles.
\draw[criticflow]
  (frozenModels.east) -| (\method.south);
\draw[supportflow]
  ([yshift=-0.35cm]frozenModels.east) -|
  ([xshift=0.82cm]\method.south);

% Classic feedback rail, separated vertically from the actor bypass.
\draw[feedback] (overviewPlant.east) -- (15.90,8.15) --
  (15.90,10.25) --
  node[pos=0.10,below,note] {$o_{t+1}$}
  (8.70,10.25) -- (observation.north);

% ===========================================================================
% (b) Deployment-time action improvement on the Figure 5 landscape.
% ===========================================================================
\fi
% Rigidly lift both lower panels: artwork, points, labels and arrows unchanged.
\ifdefined\CompactFullWidth
\begin{scope}[shift={(0,1.661111)}]
\else
\begin{scope}[shift={(0,0.75)}]
\fi
% Equal lower-panel dimensions; extra height accommodates proportional artwork.
\ifdefined\CompactFullWidth
\node[criticdetail,minimum width=7.825cm,minimum height=3.40cm]
  (landscapePanel) at (3.9125,1.16) {};
\node[paneltitle,anchor=center,text=black] at (3.9125,2.82)
  {(b) Deployment-time action improvement};
\begin{scope}[shift={(3.9125,1.16)},scale=.65,transform shape]
\begin{scope}[shift={(-3.9125,-1.16)}]
\else
\node[criticdetail,minimum width=7.825cm,minimum height=5.04cm]
  (landscapePanel) at (3.9125,1.16) {};
\fi
\ifdefined\CompactFullWidth\else
\ifdefined\HideInsetTitles\else
\node[paneltitle,anchor=center,text=black]
  at (3.9125,3.72)
  {(\ifdefined\CompactMethodPanels a\else b\fi) Deployment-time action improvement};
\fi
\fi
% The caption identifies the surface and paths as schematic, not measured iterates.
\begin{scope}
% Plot the analytic critic directly over the rectangular action domain.
\tikzset{declare function={
  criticValue(\u,\v)=1.15*exp(-((\u-0.85)^2+(\v-0.55)^2)/0.5)
    +0.42*exp(-((\u+0.9)^2+(\v+0.75)^2)/0.38)
    -0.08*((\u)^2+(\v)^2);
}}
\begin{axis}[
  % Centre the projected axis box in the space below the legend.
  at={(3.9125cm,0.90cm)},
  anchor=center,
  width=7.60cm,
  height=5.04cm,
  view={55}{37},
  domain=-1.5:1.5,
  y domain=-1.5:1.5,
  % Fine Cartesian sampling reduces visible segments along the silhouette.
  % The analytic function and rectangular domain remain unchanged.
  samples=31,
  samples y=31,
  xmin=-1.5,xmax=1.5,
  ymin=-1.5,ymax=1.5,
  zmin=-0.35,zmax=1.18,
  % Shared colour limits keep the surface and its boundary stroke matched.
  point meta min=-0.36,
  point meta max=1.07076,
  xlabel={$u_1$},
  ylabel={$u_2$},
  xlabel style={font=\sffamily\scriptsize,inner sep=1pt},
  ylabel style={font=\sffamily\scriptsize,inner sep=1pt},
  tick label style={font=\sffamily\fontsize{5.4}{6.0}\selectfont,
    text=hmInk!55},
  % Unlabelled rear-axis ticks can protrude above the surface like spikes.
  xtick=\empty,
  ytick=\empty,
  ztick=\empty,
  grid=none,
  axis line style={draw=black,line width=0.46pt},
  tick style={draw=black,line width=0.35pt},
  axis background/.style={fill=white},
  % Colour map is defined centrally in styles/paper_palette.tex.
  colormap name=referencecriticmap,
  legend style={at={(0.02,0.98)},anchor=north west,
    font=\sffamily\fontsize{5.8}{6.4}\selectfont,
    draw=none,fill=none,inner sep=1pt},
  clip=false
]
\addplot3[
  surf,
  shader=interp,
  z buffer=sort,
  opacity=0.66,
  draw=none,
  forget plot
] {criticValue(x,y)};

% Outline only the two exposed front boundaries; rear boundaries remain
% hidden by the surface. Each short stroke uses the same Q-to-colour map.
\pgfplotsextra{
  \begin{scope}[/pgf/fpu=false]
    \foreach \edge in {0,1} {
      \foreach \i in {0,...,159} {
        \pgfmathsetmacro{\edgeA}{-1.5+3*\i/160}
        \pgfmathsetmacro{\edgeB}{-1.5+3*(\i+1)/160}
        \ifnum\edge=0
          \pgfmathsetmacro{\edgeZA}{criticValue(\edgeA,-1.5)}
          \pgfmathsetmacro{\edgeZB}{criticValue(\edgeB,-1.5)}
        \else
          \pgfmathsetmacro{\edgeZA}{criticValue(1.5,\edgeA)}
          \pgfmathsetmacro{\edgeZB}{criticValue(1.5,\edgeB)}
        \fi
        \pgfmathsetmacro{\edgeMeta}{1000*((\edgeZA+\edgeZB)/2+0.36)/1.43076}
        \pgfplotscolormapaccess[0:1000]{\edgeMeta}{referencecriticmap}
        \definecolor{criticBoundaryColour}{rgb}{\pgfmathresult}
        \ifnum\edge=0
          \draw[criticBoundaryColour,opacity=0.66,line width=0.28pt,line cap=round]
            (axis cs:\edgeA,-1.5,\edgeZA) -- (axis cs:\edgeB,-1.5,\edgeZB);
        \else
          \draw[criticBoundaryColour,opacity=0.66,line width=0.28pt,line cap=round]
            (axis cs:1.5,\edgeA,\edgeZA) -- (axis cs:1.5,\edgeB,\edgeZB);
        \fi
      }
    }
  \end{scope}
}

\addplot3[
  only marks,
  mark=*,
  mark size=2.05pt,
  draw=black,
  fill=white,
  line width=0.65pt,
  forget plot
] coordinates {
  (-1.20000,-1.15000,-0.00346)
  (-1.25000,0.65000,-0.15688)
  (-0.95000,0.05000,0.00611)
  (-0.50000,1.15000,-0.11116)
  (-0.20000,-1.10000,-0.01565)
  (0.20000,-0.40000,0.07785)
  (0.50000,-1.20000,-0.13181)
  (1.20000,-0.70000,-0.11485)
  (1.30000,0.15000,0.41997)
  (1.15000,1.15000,0.25596)
  (0.20000,0.85000,0.35164)
  (0.32000,0.22000,0.51602)
};
\addplot3[
  only marks,
  mark=triangle*,
  mark size=3.2pt,
  draw=paperBC,
  fill=paperBC!20,
  forget plot
] coordinates {
  (-0.72000,-0.63000,0.29862)
  (-0.30000,0.50000,0.05672)
  (-0.40000,-0.65000,0.16812)
  (-0.80000,-1.05000,0.18345)
};

\draw[
  draw=paperCorePivot!85!black,
  line width=1.45pt,
  postaction={decorate},
  decoration={markings,mark=at position 0.5 with
    {\arrow{Stealth[length=2.3mm,width=1.6mm]}}}
]
  (axis cs:0.32000,0.22000,0.51602) --
  (axis cs:0.44000,0.30000,0.70262) --
  (axis cs:0.56000,0.37000,0.87499) --
  (axis cs:0.67000,0.44000,1.00069) --
  (axis cs:0.76000,0.49000,1.05799) --
  (axis cs:0.82136,0.53147,1.07076);
\addplot3[
  only marks,
  mark=diamond*,
  mark size=5.0pt,
  draw=overviewPolicyEdge,
  fill=overviewPolicyFill,
  line width=0.85pt,
  forget plot
] coordinates {(0.32000,0.22000,0.51602)};
% Keep the original uniform-seed circle visible inside its selection diamond.
\addplot3[
  only marks,
  mark=*,
  mark size=2.05pt,
  draw=black,
  fill=white,
  line width=0.65pt,
  forget plot
] coordinates {(0.32000,0.22000,0.51602)};
\addplot3[
  only marks,
  mark=star,
  mark size=3.2pt,
  draw=paperCorePivot!85!black,
  fill=paperCorePivot!85!black,
  forget plot
] coordinates {(0.82136,0.53147,1.07076)};
\node[annotation,anchor=east,font=\sffamily\fontsize{5.8}{6.5}\selectfont,
  text=paperBC,fill=none,text opacity=1]
  at ([xshift=-4pt]axis cs:-0.72,-0.63,0.29862) {BC sample};
\node[annotation,anchor=west,font=\sffamily\fontsize{5.8}{6.5}\selectfont,
  text=overviewPolicyEdge,fill=none,text opacity=1]
  at ([xshift=14pt,yshift=8pt]axis cs:0.32,0.22,0.51602) {highest-ranked seed};

\end{axis}
\end{scope}
\ifdefined\CompactFullWidth\else
\node[annotation,anchor=west,font=\sffamily\fontsize{6.0}{7.0}\selectfont,
  text=hmInk] at (0.65,3.28)
  {$\circ$\; other candidates};
\node[annotation,anchor=west,font=\sffamily\fontsize{6.0}{7.0}\selectfont,
  text=paperBC] at (2.75,3.28)
  {$\triangle$\; BC sample};
\node[annotation,anchor=west,font=\sffamily\fontsize{6.0}{7.0}\selectfont,
  text=paperCorePivot] at (4.95,3.28)
  {$\rightarrow$\; selected refinement};
\fi
\node[annotation,font=\sffamily\fontsize{5.7}{6.2}\selectfont,
  text=hmInk!72,rotate=90]
  at (0.25,0.95) {$\widehat Q(o_t,u)$};
\ifdefined\CompactFullWidth\else
\node[annotation,font=\sffamily\fontsize{6.4}{7.4}\selectfont,
  text=hmInk!85,align=center] at (3.9125,-1.30)
  {Rank seeds by $Q$; refine only the highest-ranked seed};
\fi

% ===========================================================================
% (c) Physical closed-loop control context.
% ===========================================================================
\ifdefined\CompactFullWidth
\end{scope}
\end{scope}
\node[annotation,anchor=west,font=\sffamily\fontsize{6.0}{7.0}\selectfont,
  text=hmInk] at (.45,2.38) {$\circ$\; other candidates};
\node[annotation,anchor=west,font=\sffamily\fontsize{6.0}{7.0}\selectfont,
  text=paperBC] at (2.65,2.38) {$\triangle$\; BC sample};
\node[annotation,anchor=west,font=\sffamily\fontsize{6.0}{7.0}\selectfont,
  text=paperCorePivot] at (4.85,2.38) {$\rightarrow$\; selected refinement};
\node[annotation,font=\sffamily\fontsize{6.4}{7.4}\selectfont,
  text=hmInk!85,align=center] at (3.9125,-.44)
  {Rank seeds by $Q$; refine only the highest-ranked seed};
\node[taskdetail,minimum width=7.825cm,minimum height=3.40cm]
  (taskPanel) at (12.0875,1.16) {};
\node[paneltitle,anchor=center,text=black] at (12.0875,2.82)
  {(c) Drone control task};
\begin{scope}[shift={(12.0875,1.16)},scale=.65,transform shape]
\begin{scope}[shift={(-12.0875,-1.16)}]
\else
\node[taskdetail,minimum width=7.825cm,minimum height=5.04cm]
  (taskPanel) at (12.0875,1.16) {};
\fi
\ifdefined\CompactFullWidth\else
\ifdefined\HideInsetTitles\else
\node[paneltitle,anchor=center,text=black]
  at (12.0875,3.72) {(\ifdefined\CompactMethodPanels b\else c\fi) Drone control task};
\fi
\fi

% Compact adaptation of Figure 3: projected tracking geometry, smoothed
% waypoint path, persistent side wind, and the detailed racing quadrotor.
% Uniform scaling preserves the reference geometry and drone aspect ratio.
% Centre the complete projected illustration, including its axis labels.
\begin{scope}[shift={(6.15,-1.14)},scale=0.90,transform shape]
  % Projected coordinate frame and ground-plane waypoint projection.
  \fill[overviewDPFill!35] (2.80,0.35) -- (9.00,0.35) -- (9.75,1.60) --
    (3.55,1.60) -- cycle;
  \foreach \x in {2.80,3.80,...,8.80} {
    \draw[very thin,overviewDPEdge!7] (\x,0.35) -- ++(0.75,1.25);
  }
  \foreach \s in {0.00,0.25,0.50,0.75,1.00} {
    \draw[very thin,overviewDPEdge!7]
      ({2.80+0.75*\s},{0.35+1.25*\s}) --
      ({9.00+0.75*\s},{0.35+1.25*\s});
  }
  \draw[-{Stealth[length=5pt,width=5.4pt]},overviewDPEdge!46,line width=0.62pt]
    (2.80,0.35) -- (10.10,0.35)
    node[right,font=\sffamily\footnotesize,text=black] {$x$};
  \draw[-{Stealth[length=5pt,width=5.4pt]},overviewDPEdge!46,line width=0.62pt]
    (2.80,0.35) -- (3.50,1.55)
    node[above right,font=\sffamily\footnotesize,text=black] {$y$};
  \draw[-{Stealth[length=5pt,width=5.4pt]},overviewDPEdge!46,line width=0.62pt]
    (2.80,0.35) -- (2.80,4.65)
    node[above,font=\sffamily\footnotesize,text=black] {$z$};

  % Expanded flight path, with p1 and p2 advanced slightly farther.
  % Ground projections follow the horizontal waypoint positions.
  \coordinate (taskGroundOne) at (5.70,0.72);
  \coordinate (taskGroundTwo) at (9.45,1.10);
  \coordinate (taskGroundThree) at (6.6125,1.60);
  \coordinate (taskGroundFour) at (3.30,1.08);
  \coordinate (taskPointOne) at (5.70,1.40);
  \coordinate (taskPointTwo) at (9.45,2.60);
  \coordinate (taskPointThree) at (6.6125,4.0875);
  \coordinate (taskPointFour) at (3.30,2.3375);
  \coordinate (taskTwoIn) at ($(taskPointOne)!0.78!(taskPointTwo)$);
  \coordinate (taskTwoOut) at ($(taskPointTwo)!0.22!(taskPointThree)$);
  \coordinate (taskThreeIn) at ($(taskPointTwo)!0.78!(taskPointThree)$);
  \coordinate (taskThreeOut) at ($(taskPointThree)!0.22!(taskPointFour)$);

  % Ground projection and the smoothed nominal reference from Figure 3.
  \draw[thick,overviewActionEdge,densely dashed,opacity=0.48]
    (taskGroundOne) -- (taskGroundTwo) --
    (taskGroundThree) -- (taskGroundFour);
  \foreach \ground/\flight in {
    taskGroundOne/taskPointOne,taskGroundTwo/taskPointTwo,
    taskGroundThree/taskPointThree,taskGroundFour/taskPointFour} {
    \draw[densely dotted,overviewDPEdge!38,line width=0.58pt]
      (\ground) -- (\flight);
  }
  \draw[very thick,paperReference,densely dashed,line cap=round,line join=round]
    (taskPointOne) -- (taskTwoIn)
    .. controls (taskPointTwo) .. (taskTwoOut) -- (taskThreeIn)
    .. controls (taskPointThree) .. (taskThreeOut) -- (taskPointFour);
  \draw[->,thick,paperReference,shorten >=1pt]
    ($(taskPointOne)!0.42!(taskTwoIn)$) --
    ($(taskPointOne)!0.68!(taskTwoIn)$);
  \draw[->,thick,paperReference,shorten >=1pt]
    ($(taskTwoOut)!0.38!(taskThreeIn)$) --
    ($(taskTwoOut)!0.64!(taskThreeIn)$);
  \draw[->,thick,paperReference,shorten >=1pt]
    ($(taskThreeOut)!0.35!(taskPointFour)$) --
    ($(taskThreeOut)!0.62!(taskPointFour)$);

  \node[annotation,font=\sffamily\footnotesize,text=overviewDPEdge!62]
    at ($(taskPointOne)+(0.12,-0.16)$) {$p_1$};
  \node[annotation,font=\sffamily\footnotesize,text=overviewActionEdge]
    at ($(taskPointTwo)+(0.14,-0.15)$) {$p_2$};
  \node[annotation,font=\sffamily\footnotesize,text=overviewActionEdge]
    at ($(taskPointThree)+(0.00,0.18)$) {$p_3$};
  \node[annotation,font=\sffamily\footnotesize,text=overviewActionEdge]
    at ($(taskPointFour)+(-0.27,0.02)$) {$p_4$};
  % Persistent side-wind field, simplified from Figure 3.
  \fill[paperPropellerLight,draw=paperPropellerEdge,line width=0.38pt,dashed,
    fill opacity=0.48,draw opacity=0.85]
    (3.60,1.31) .. controls (4.48,1.18) and (5.38,1.28) .. (6.25,1.46)
    .. controls (6.38,1.53) and (6.38,1.75) .. (6.24,1.92)
    .. controls (5.36,1.88) and (4.46,1.88) .. (3.58,2.00)
    .. controls (3.42,1.78) and (3.42,1.52) .. cycle;
  \draw[-{Stealth[length=8pt,width=8pt]},paperPropellerEdge,densely dashed,
    line width=1.10pt,
    line cap=round,line join=round]
    (3.68,1.38) .. controls (4.55,1.28) and (5.35,1.36) .. (5.84,1.47);
  \draw[-{Stealth[length=8pt,width=8pt]},paperPropellerEdge,densely dashed,
    line width=1.10pt,
    line cap=round,line join=round]
    (3.62,1.88) .. controls (4.52,1.76) and (5.34,1.78) .. (5.84,1.82);
  \node[annotation,font=\sffamily\small,text=paperPropellerEdge,anchor=west]
    at (3.58,1.02) {side wind $w_t$};

  % Racing quadrotor and its ground projection.
  \coordinate (taskProjection) at ($(taskGroundOne)!0.42!(taskGroundTwo)$);
  \coordinate (taskDrone) at ($(taskPointOne)!0.42!(taskPointTwo)$);
  \fill[overviewDPEdge!25] (taskProjection)
    ellipse[x radius=0.42cm,y radius=0.09cm];
  \draw[dashed,overviewDPEdge!38,line width=0.72pt]
    (taskDrone) -- (taskProjection);
  \begin{scope}[shift={(taskDrone)},rotate=16,xscale=1.10,yscale=0.86]
    % Propeller-wash halos from Figure 3.
    \foreach \x/\y/\tilt in {-0.78/-0.54/-4,0.78/0.54/-4,
      -0.78/0.54/4,0.78/-0.54/4} {
      \begin{scope}[shift={(\x,\y)},rotate=\tilt]
        \draw[draw=paperPropellerEdge,fill=paperPropellerLight,
          line width=0.42pt,opacity=0.72]
          (0,0) ellipse[x radius=0.34,y radius=0.19];
      \end{scope}
    }

    % Silver frame, motor pods, and fuselage; method colour stays in the logic.
    \draw[line width=2.8pt,droneLight!48,line cap=round]
      (-0.78,-0.54) -- (0.78,0.54);
    \draw[line width=2.8pt,droneLight!48,line cap=round]
      (-0.78,0.54) -- (0.78,-0.54);
    \draw[line width=1.65pt,droneAccent,line cap=round]
      (-0.78,-0.54) -- (0.78,0.54);
    \draw[line width=1.65pt,droneAccent,line cap=round]
      (-0.78,0.54) -- (0.78,-0.54);
    \foreach \x/\y in {-0.78/-0.54,0.78/0.54,-0.78/0.54,0.78/-0.54} {
      \shade[inner color=droneLight!55!white,
        outer color=droneAccent,draw=droneAccent!85!black,
        line width=0.32pt] (\x,\y) circle[radius=0.17];
      \draw[fill=droneAccent!82!black,draw=black,line width=0.22pt]
        (\x,\y) circle[radius=0.095];
    }
    \shade[left color=droneAccent!88!black,right color=droneAccent,
      middle color=droneLight,draw=black,
      line width=0.52pt,rounded corners=1.2pt]
      (-0.48,-0.19) -- (0.18,-0.22) -- (0.55,0) -- (0.18,0.22) --
      (-0.48,0.19) -- (-0.62,0) -- cycle;
    \draw[fill=droneAccent,draw=black,line width=0.30pt]
      (0.46,0) -- (0.66,0.11) -- (0.66,-0.11) -- cycle;

    % Spinning three-blade propellers and motor caps.
    \foreach \x/\y/\spin in {-0.78/-0.54/18,0.78/0.54/18,
      -0.78/0.54/-18,0.78/-0.54/-18} {
      \begin{scope}[shift={(\x,\y)}]
        \foreach \ang in {0,120,240} {
          \pgfmathsetmacro{\bladeangle}{\ang+\spin}
          \begin{scope}[rotate=\bladeangle]
            \shade[left color=paperPropellerLight!45!white,
              right color=paperPropellerLight,draw=paperPropellerEdge,
              line width=0.22pt,opacity=0.78]
              (0.045,0.017) .. controls (0.19,0.070) and
              (0.32,0.060) .. (0.40,0) .. controls
              (0.31,-0.055) and (0.18,-0.050) ..
              (0.045,-0.017) -- cycle;
          \end{scope}
        }
        \shade[inner color=droneLight!55!white,
          outer color=droneAccent,draw=droneAccent!85!black,
          line width=0.30pt] (0,0) circle[radius=0.125];
        \draw[fill=droneAccent,draw=black,line width=0.18pt]
          (0,0) circle[radius=0.060];
      \end{scope}
    }

    % Camera lens and navigation LEDs.
    \shade[inner color=paperPropellerLight,outer color=paperPropellerEdge,
      draw=black,line width=0.25pt] (0.61,0) circle[radius=0.066];
    \fill[white,opacity=0.82] (0.63,0.020) circle[radius=0.014];
    \draw[fill=white,draw=black,line width=0.20pt]
      (-0.49,0.13) circle[radius=0.024];
    \draw[fill=white,draw=black,line width=0.20pt]
      (-0.49,-0.13) circle[radius=0.024];
    \draw[fill=paperPropellerEdge,draw=black,line width=0.18pt]
      (0.31,0.15) circle[radius=0.021];
    \draw[fill=paperPropellerEdge,draw=black,line width=0.18pt]
      (0.31,-0.15) circle[radius=0.021];
  \end{scope}

\end{scope}

\ifdefined\CompactFullWidth
\end{scope}
\end{scope}
\fi
\end{scope}

% Exclude invisible plotting margins when sizing the complete illustration.
\pgfresetboundingbox
\ifdefined\CompactMethodPanels
\path[use as bounding box] (-0.04,-0.93) rectangle (16.04,4.85);
\else
\ifdefined\CompactFullWidth
\path[use as bounding box] (-0.04cm,0.78cm) rectangle (16.04cm,10.43cm);
\else
\path[use as bounding box] (-0.04,-0.93) rectangle (16.04,11.54);
\fi
\fi
\end{tikzpicture}
\endgroup
